\begin{ExerciseBlock}
\ExerciseHeading{}{Verständnisfragen}{work_and_energy_checks}
Entscheiden Sie, welche Aussagen wahr oder falsch sind, und begründen Sie Ihre Entscheidung.
\begin{enumerate}
\item Nur Kräfte senkrecht zur zurückgelegten Strecke, tragen zur physikalischen Arbeit bei.\TrueFalseAnswer{falsch}
\item Bei einer konstanten Kraft in einer Dimension ist die Änderung der potenziellen Energie proportional zur Ortsänderung.\TrueFalseAnswer{wahr}
\item Die potenzielle Energie einer idealen Feder ist proportional zum Quadrat der Auslenkung aus der Ruhelage.\TrueFalseAnswer{wahr}
\item Verdoppelt sich die Geschwindigkeit eines Objektes, so verdoppelt sich auch seine kinetische Energie.\TrueFalseAnswer{falsch}
\item Die physikalische Arbeit ergibt sich als Fläche unter dem $F_{||}$-$s$-Diagramm.\TrueFalseAnswer{wahr}
\item Die Gesamtenergie eines abgeschlossenen Systems ist immer Null.\TrueFalseAnswer{falsch}
\item Potentielle Energie $E_{pot}(x)$ ist die negative Stammfunktion der dazugehörigen Kraftfunktion $F(x)$.\TrueFalseAnswer{wahr}
\item Die Methode der Energieerhaltung erlaubt die Richtung einer Geschwindigkeit zu bestimmen.\TrueFalseAnswer{falsch}
\item Leistung entspricht die Fläche unter der Arbeitsfunktion $W(t)$.\TrueFalseAnswer{falsch}
\item Jede reale Maschine hat einen Wirkungsgrad $\eta < 1$.\TrueFalseAnswer{wahr}
\end{enumerate}
\end{ExerciseBlock}
